\documentclass[UTF8, 12pt, fontset=fandol, oneside]{ctexart}
\usepackage{researchnotes}

\title{第四章\quad 线性映射}
\author{}
\date{}

\begin{document}

\maketitle

\section*{§4.1\quad 线性映射的概念}

读者已经学过映射的概念，我们现在来复习一下. 所谓映射，是指从一个集合 $A$ 到另一个集合 $B$ 的对应 $\varphi:A\to B$. 对 $A$ 中任一元素 $a$，均有唯一的元素 $b\in B$ 与之对应，记为 $b=\varphi(a)$. 元素 $b$ 称为 $a$ 在 $\varphi$ 下的像，$a$ 称为元素 $b$ 的原像或逆像. $A$ 中元素在 $\varphi$ 下的像全体构成 $B$ 的一个子集，记为 $\varphi(A)$ 或 $\operatorname{Im}\varphi$. 如果 $\operatorname{Im}\varphi=B$，即对 $B$ 中任一元素 $b$，在 $A$ 中均有元素 $a$，使 $b=\varphi(a)$，则称 $\varphi$ 是满映射或称 $\varphi$ 是映上的映射. 如果映射 $\varphi$ 适合下列条件：若 $a\neq a'$，则 $\varphi(a)\neq\varphi(a')$，那么就称 $\varphi$ 是单映射. 单映射的另外一个等价说法是从 $\varphi(a)=\varphi(a')$ 可推出 $a=a'$. 如果 $\varphi$ 既是单映射又是满映射，则称 $\varphi$ 是双射. 这时不仅对 $A$ 中的任一元素，有且仅有 $B$ 中的一个元素与之对应，而且对 $B$ 中的任一元素，有且仅有 $A$ 中的一个元素与之对应. 因此，双射又称为一一对应.

函数 $y=x^2$ 可以看成是实数集 $\mathbb{R}$ 到 $\mathbb{R}$ 的映射. 这个映射不是单映射也不是满映射. 因为 $-1\neq 1$，但 $(-1)^2=1^2$. 另一方面，$\mathbb{R}$ 中的元素 $-1$ 没有原像. 这个函数的像为 $[0,+\infty)$. 如果我们把 $y=x^2$ 看成是 $(-\infty,+\infty)\to[0,+\infty)$ 的映射，那么这是个映上的映射，但仍不是单映射.

设 $M_n(\mathbb{R})$ 是实数域 $\mathbb{R}$ 上的 $n(n>1)$ 阶方阵全体组成的集合，定义 $M_n(\mathbb{R})\to\mathbb{R}$ 的映射为：$\varphi(A)=\det A$，则 $\varphi$ 是映射，且这是个映上的映射，但它不是单映射. 事实上，对 $\mathbb{R}$ 中任意一个实数，均有无穷多个矩阵，其行列式的值等于这个实数.

设 $y=x^3$ 是 $[0,1)$ 到 $[0,+\infty)$ 上的函数，这显然是一个映射且是一个单映射，但它不是满映射.

Descartes 平面上的点到实数对之间的对应：
\[
\varphi:C\mapsto(a,b),
\]
其中点 $C$ 的横坐标为 $a$，纵坐标为 $b$. 这是一个映射且是一个双射，即一一对应.

若集合 $A$ 是集合 $B$ 的子集，作 $A\to B$ 的映射 $j$：
\[
j(a)=a,\quad a\in A,
\]
则 $j$ 是一个映射且显然是单映射. 若 $A=B$，则 $j$ 是一个一一对应，这时映射 $j$ 实际上把 $A$ 中任一元素映射为自身，因此称为恒等映射，记为 $1_A$ 或 $I_A$.

一个集合 $A$ 到自身的映射通常称为变换，比如 $y=x^2$ 可以看成是 $\mathbb{R}$ 自身的变换.

集合 $A$ 到 $B$ 的两个映射 $f$ 与 $g$，若对任意的 $a\in A$，都有 $f(a)=g(a)$，则称它们相等，记为 $f=g$.

若 $f$ 是集合 $A\to B$ 的映射，$g$ 是集合 $B\to C$ 的映射，定义映射 $g$ 与 $f$ 的复合 $g\circ f$ 为集合 $A\to C$ 的映射，且
\[
(g\circ f)(a)=g(f(a)),\quad a\in A.
\]
这里需要注意的是并非任意两个映射都有复合，只有当 $f$ 的像落在 $g$ 所定义的集合上时才可定义 $g$ 与 $f$ 的复合. 如果 $A=B$，那么 $A$ 上的任意两个变换都可复合.

若 $f$ 是 $A\to B$ 的映射，$g$ 是 $B\to C$ 的映射，$h$ 是 $C\to D$ 的映射，则 $h\circ(g\circ f)$ 及 $(h\circ g)\circ f$ 都是 $A\to D$ 的映射且对任意的 $a\in A$，有
\[
((h\circ g)\circ f)(a)=(h\circ g)(f(a))=h(g(f(a))),
\]
\[
(h\circ(g\circ f))(a)=h((g\circ f)(a))=h(g(f(a))),
\]
因此
\begin{equation}
(h\circ g)\circ f=h\circ(g\circ f).
\tag{4.1.1}
\end{equation}
上式通常称为映射复合的结合律. 正因为如此，我们写 3 个 (或 3 个以上) 映射的复合时常常略去括号，写为 $h\circ g\circ f$. 通常复合号“$\circ$”也省略，即 $g\circ f$ 写为 $gf$.

下面我们着重讨论一下双射. 设 $f$ 是 $A\to B$ 的双射，我们定义 $B\to A$ 的映射 $g$ 如下：对任一 $b\in B$，取 $b$ 在 $f$ 下的原像记为 $a$，定义 $g(b)=a$. 由于 $f$ 是双射，故对 $B$ 中的元素 $b$，有且仅有一个 $a$ 作为 $b$ 在 $f$ 下的原像，因此 $g$ 确是 $B\to A$ 的映射. 不仅如此，显然 $g$ 也是一个双射，且
\[
gf=1_A,\quad fg=1_B.
\]
我们称 $g$ 是 $f$ 的逆映射，记为 $g=f^{-1}$.

\noindent\textbf{命题 4.1.1}\quad 设 $f$ 是集合 $A\to B$ 的映射，如果存在 $B\to A$ 的映射 $g$，使
\[
gf=1_A,\quad fg=1_B,
\]
则 $f$ 是双射且 $g=f^{-1}$.

\noindent\textbf{证明}\quad 先证 $f$ 是单映射：若 $f(a_1)=f(a_2)$，则由 $gf(a_1)=1_A(a_1)=a_1$，$gf(a_2)=1_A(a_2)=a_2$，知 $a_1=a_2$，因此 $f$ 为单映射.

又对 $B$ 中任一元素 $b$，$g(b)\in A$ 且 $fg(b)=1_B(b)=b$. 因此 $g(b)$ 是 $b$ 在 $f$ 下的原像，即 $f$ 是映上的映射. $\square$

从命题 4.1.1 的证明中可看出，从 $gf=1_A$ 可推出 $f$ 是单映射，从 $fg=1_B$ 可推出 $f$ 是满映射.

现在我们转而来考虑线性映射.

\noindent\textbf{定义 4.1.1}\quad 设 $\varphi$ 是数域 $\mathbb{K}$ 上线性空间 $V$ 到 $\mathbb{K}$ 上线性空间 $U$ 的映射，如果 $\varphi$ 适合下列条件：
\[
(1)\ \varphi(\alpha+\beta)=\varphi(\alpha)+\varphi(\beta),\quad \alpha,\beta\in V;
\]
\[
(2)\ \varphi(k\alpha)=k\varphi(\alpha),\quad k\in\mathbb{K},\ \alpha\in V,
\]
则称 $\varphi$ 是 $V\to U$ 的线性映射. $V$ 到自身的线性映射称为 $V$ 上的线性变换. 若 $\varphi$ 作为映射是单的，则称 $\varphi$ 是单线性映射；若 $\varphi$ 作为映射是满的，则称 $\varphi$ 是满线性映射. 若 $\varphi$ 是双射，则称 $\varphi$ 是线性同构，简称同构. 若 $V=U$，则 $V$ 自身上的同构称为自同构.

请读者注意，我们曾在第三章中定义了两个线性空间的同构，那个定义与现在的定义是相同的.

\noindent\textbf{例 4.1.1}\quad 设 $V,U$ 是数域 $\mathbb{K}$ 上的线性空间，定义 $\varphi$ 为 $V\to U$ 的映射：对任意的 $\alpha\in V$，$\varphi(\alpha)=0$. 显然 $\varphi$ 是一个线性映射，称为零映射，通常记为 0，但要注意这是一个映射.

\noindent\textbf{例 4.1.2}\quad 设 $V$ 是 $\mathbb{K}$ 上的线性空间，显然 $V$ 到自身的恒等映射 $1_V$ 是 $V$ 上的线性变换，称为恒等变换，记为 $I_V$ 或 $\operatorname{Id}_V$，在不至于混淆的情形下，也简记为 $I$.

\noindent\textbf{例 4.1.3}\quad 设 $V=\mathbb{K}_n$，$U=\mathbb{K}_m$ 分别是数域 $\mathbb{K}$ 上的 $n$ 维和 $m$ 维列向量空间，$A$ 是 $\mathbb{K}$ 上的 $m\times n$ 矩阵，定义 $V\to U$ 的映射 $\varphi$ 为
\[
\varphi(\alpha)=A\alpha.
\]
这个映射由矩阵乘法定义 ($m\times n$ 矩阵乘以 $n$ 维列向量是一个 $m$ 维列向量)，由矩阵乘法性质容易验证 $\varphi$ 是一个线性映射.

\noindent\textbf{例 4.1.4}\quad 设 $\mathbb{K}^n$ 是 $\mathbb{K}$ 上 $n$ 维行向量空间，$\mathbb{K}_n$ 是 $\mathbb{K}$ 上 $n$ 维列向量空间，定义 $\mathbb{K}^n\to\mathbb{K}_n$ 的映射 $\varphi$：
\[
(a_1,a_2,\cdots,a_n)\mapsto
\left(\begin{array}{c}
a_1\\
a_2\\
\vdots\\
a_n
\end{array}\right),
\]
容易验证 $\varphi$ 是线性同构.

\noindent\textbf{例 4.1.5}\quad 设 $V$ 是 $\mathbb{K}$ 上的线性空间，定义 $V$ 上的变换 $\varphi$：对任意的 $\alpha\in V$，有
\[
\varphi(\alpha)=k\alpha,
\]
其中 $k$ 是一个固定常数. 容易验证 $\varphi$ 是 $V$ 上的线性变换，这个变换常称为纯量变换或数量变换.

\noindent\textbf{例 4.1.6}\quad 设 $V$ 是区间 $[0,1]$ 上的实无穷次可微函数全体组成的实线性空间，定义 $\varphi$ 为求导变换
\[
\varphi(f(x))=\frac{\mathrm{d}}{\mathrm{d}x}f(x),
\]
由求导性质知 $\varphi$ 是 $V$ 上的线性变换.

\noindent\textbf{命题 4.1.2}\quad 设 $\varphi$ 是 $V\to U$ 的线性映射，则

(1) $\varphi(0)=0$;

(2) $\varphi(k\alpha+l\beta)=k\varphi(\alpha)+l\varphi(\beta)$，$\alpha,\beta\in V$，$k,l\in\mathbb{K}$;

(3) 若 $\varphi$ 是同构，则其逆映射 $\varphi^{-1}$ 也是线性映射，从而是 $U\to V$ 的同构.

\noindent\textbf{证明}\quad (1) $\varphi(0)=\varphi(0\cdot\alpha)=0\cdot\varphi(\alpha)=0$.

(2) $\varphi(k\alpha+l\beta)=\varphi(k\alpha)+\varphi(l\beta)=k\varphi(\alpha)+l\varphi(\beta)$.

(3) 见第三章定理 3.7.2 的证明. $\square$

\noindent\textbf{注}\quad 若 $V\to U$ 的一个映射 $\varphi$ 适合命题 4.1.2 中的 (2)，则 $\varphi$ 必是线性映射. 读者不难自己验证这一点.

最后需要提醒读者注意的是，在线性映射的定义中，要求 $V$ 与 $U$ 都是数域 $\mathbb{K}$ 上的线性空间，不同数域上线性空间之间的映射不是线性映射.

\subsection*{习题 4.1}

1. 判断下列映射是否是线性变换：

(1) 设 $V$ 是 Descartes 平面，$\varphi$ 把平面上任一向量伸长 $n$ 倍 ($n$ 是固定的自然数)；

(2) 设 $V$ 是 Descartes 平面，$\varphi$ 把平面上任一向量逆时针转动 $60^\circ$，但其长度保持不变；

(3) 设 $V$ 是 $[0,1]$ 区间上连续函数全体组成的实线性空间，$\varphi$ 是 $V$ 上的变换：对任意的 $f(x)\in V$，
\[
\varphi(f(x))=\int_0^x f(t)\mathrm{d}t;
\]

(4) 设 $V$ 是 Descartes 平面，$\varphi$ 为 $V$ 上的变换：
\[
\varphi(x,y)=(2x^2,y),\quad (x,y)\in V;
\]

(5) 设 $V$ 是 Descartes 平面，$(a,b)$ 是平面上固定的一点，$\varphi$ 是 $V$ 上的变换：
\[
\varphi(x,y)=(x+a,y+b),\quad (x,y)\in V.
\]

2. 证明：$\mathbb{K}_n$ 上的变换
\[
\varphi(\alpha)=A\alpha+\beta
\]
是线性变换的充分必要条件是 $\beta=0$. 这里 $A$ 是一个 $n$ 阶矩阵，$\beta$ 是固定的 $n$ 维列向量，$\alpha$ 是 $\mathbb{K}_n$ 中任一向量.

3. 设线性空间 $V=V_1\oplus V_2$，并且 $\varphi_1$ 及 $\varphi_2$ 分别是 $V_1,V_2$ 到 $U$ 的线性映射，证明：存在唯一的从 $V$ 到 $U$ 的线性映射 $\varphi$，当 $\varphi$ 限制在 $V_i$ 上时等于 $\varphi_i(i=1,2)$.

4. 设 $V$ 是由几乎处处为零的无穷实数数列 (即 $(a_0,a_1,a_2,\cdots,a_n,\cdots)$，其中只有有限多个 $a_i$ 不为零) 组成的实向量空间，$\mathbb{R}[x]$ 是所有实系数多项式组成的实向量空间. 定义 $\varphi$ 如下：
\[
\varphi(a_0,a_1,a_2,\cdots,a_n,\cdots)=a_0+a_1x+a_2x^2+\cdots+a_nx^n,
\]
其中 $a_n\neq 0$，而 $a_s=0(s>n)$. 求证：$\varphi$ 是线性同构.

\section*{§4.2\quad 线性映射的运算}

在例 4.1.3 中，我们已经知道任意一个 $m\times n$ 矩阵 $A$ 都可以定义一个从 $n$ 维列向量空间到 $m$ 维列向量空间的线性映射：$\varphi(\alpha)=A\alpha$. 如果另有一个 $m\times n$ 矩阵 $B$，它定义的线性映射是 $\psi(\alpha)=B\alpha$，注意到矩阵之间存在的运算，我们可以定义这两个映射的加法：$(\varphi+\psi)(\alpha)=(A+B)\alpha$. 显然这仍是一个线性映射. 类似地，我们还可定义线性映射的数乘：$(k\varphi)(\alpha)=kA\alpha$. 对一般的线性映射，我们是否也可以定义它们的加法和数乘呢？

\noindent\textbf{定义 4.2.1}\quad 设 $\varphi,\psi$ 是 $\mathbb{K}$ 上线性空间 $V\to U$ 的线性映射，定义 $\varphi+\psi$ 为 $V\to U$ 的映射：
\[
(\varphi+\psi)(\alpha)=\varphi(\alpha)+\psi(\alpha),\quad \alpha\in V.
\]
若 $k\in\mathbb{K}$，定义 $k\varphi$ 为 $V\to U$ 的映射：
\[
(k\varphi)(\alpha)=k\varphi(\alpha),\quad \alpha\in V.
\]

容易验证 $\varphi+\psi$ 是线性映射. 事实上
\[
\begin{aligned}
(\varphi+\psi)(k\alpha+l\beta)
&=\varphi(k\alpha+l\beta)+\psi(k\alpha+l\beta)\\
&=\varphi(k\alpha)+\varphi(l\beta)+\psi(k\alpha)+\psi(l\beta)\\
&=k\varphi(\alpha)+l\varphi(\beta)+k\psi(\alpha)+l\psi(\beta)\\
&=k(\varphi(\alpha)+\psi(\alpha))+l(\varphi(\beta)+\psi(\beta))\\
&=k(\varphi+\psi)(\alpha)+l(\varphi+\psi)(\beta).
\end{aligned}
\]
同理可证明 $k\varphi$ 也是线性映射.

现在我们考虑从 $V$ 到 $U$ 的线性映射全体组成的集合 $\mathcal{L}(V,U)$. 在这个集合上，既然我们定义了加法和数乘，那它是一个线性空间吗？

\noindent\textbf{命题 4.2.1}\quad 设 $\mathcal{L}(V,U)$ 是 $V\to U$ 的线性映射全体，则在上述线性映射的加法及数乘定义下，$\mathcal{L}(V,U)$ 是 $\mathbb{K}$ 上的线性空间. 特别，$V\to\mathbb{K}$ 的所有线性函数全体构成一个线性空间.

这个命题的证明很容易，只需按照线性空间的定义逐条验证即可，我们将证明留给读者.

若 $U=\mathbb{K}$，即把 $\mathbb{K}$ 看成是 $\mathbb{K}$ 上的一维空间，则 $V\to\mathbb{K}$ 的线性映射通常称为 $V$ 上的线性函数. $V$ 上所有的线性函数构成的线性空间通常称为 $V$ 的共轭空间，记为 $V^*$. 当 $V$ 是有限维空间时，$V^*$ 也称为 $V$ 的对偶空间.

若 $V=U$，我们用 $\mathcal{L}(V)$ 来记 $\mathcal{L}(V,V)$，即 $V$ 上线性变换全体构成的线性空间. 这时在 $\mathcal{L}(V)$ 上，除了加法和数乘运算外，还有乘法运算，这个乘法就是映射的复合.

\noindent\textbf{定义 4.2.2}\quad 设 $A$ 是数域 $\mathbb{K}$ 上的线性空间，如果在 $A$ 上定义了一个乘法“$\cdot$” (通常可以省略)，使对任意的 $A$ 中元素 $a,b,c$ 及 $\mathbb{K}$ 中元素 $k$，适合下列条件：

(1) 乘法结合律：$a\cdot(b\cdot c)=(a\cdot b)\cdot c$;

(2) 存在 $A$ 中元 $e$，使对一切 $a\in A$，均有
\[
e\cdot a=a\cdot e=a;
\]

(3) 分配律：
\[
a\cdot(b+c)=a\cdot b+a\cdot c,
\]
\[
(b+c)\cdot a=b\cdot a+c\cdot a;
\]

(4) 乘法与数乘的相容性：
\[
(ka)\cdot b=k(a\cdot b)=a\cdot(kb),
\]
则称 $A$ 是数域 $\mathbb{K}$ 上的代数，元素 $e$ 称为 $A$ 的恒等元.

\noindent\textbf{注}\quad $A$ 的恒等元常常用 1 表示，注意不要与数 1 混淆.

\noindent\textbf{定理 4.2.1}\quad 设 $V$ 是数域 $\mathbb{K}$ 上的线性空间，则 $\mathcal{L}(V)$ 是 $\mathbb{K}$ 上的代数.

\noindent\textbf{证明}\quad 由命题 4.2.1，$\mathcal{L}(V)$ 是 $\mathbb{K}$ 上的线性空间. 我们逐条来验证定义 4.2.2.

(1) 乘法结合律实际上就是映射复合的结合律，因此自然成立.

(2) 设 $1_V$ 是 $V$ 上的恒等映射，由上节可知它是线性变换，显然对任意的 $\varphi\in\mathcal{L}(V)$，有
\[
1_V\circ\varphi=\varphi\circ 1_V=\varphi,
\]
因此 $1_V$ 是 $\mathcal{L}(V)$ 的恒等元.

(3) 设 $\varphi_1,\varphi_2,\varphi_3$ 都是 $V$ 上线性变换，对任意的 $\alpha\in V$，有
\[
\begin{aligned}
(\varphi_1\circ(\varphi_2+\varphi_3))(\alpha)
&=\varphi_1((\varphi_2+\varphi_3)(\alpha))\\
&=\varphi_1(\varphi_2(\alpha)+\varphi_3(\alpha))\\
&=\varphi_1(\varphi_2(\alpha))+\varphi_1(\varphi_3(\alpha))\\
&=(\varphi_1\circ\varphi_2)(\alpha)+(\varphi_1\circ\varphi_3)(\alpha)\\
&=(\varphi_1\circ\varphi_2+\varphi_1\circ\varphi_3)(\alpha).
\end{aligned}
\]
因此
\[
\varphi_1\circ(\varphi_2+\varphi_3)=\varphi_1\circ\varphi_2+\varphi_1\circ\varphi_3.
\]
同理可证明另外一个分配律.

(4) 设 $k$ 是 $\mathbb{K}$ 中任一数，$\varphi,\psi$ 为 $V$ 上线性变换，则对任意的 $\alpha\in V$，有
\[
\begin{aligned}
((k\varphi)\circ\psi)(\alpha)
&=(k\varphi)(\psi(\alpha))=k(\varphi(\psi(\alpha)))\\
&=k((\varphi\circ\psi)(\alpha))=(k(\varphi\circ\psi))(\alpha),
\end{aligned}
\]
从而
\[
(k\varphi)\circ\psi=k(\varphi\circ\psi).
\]
同理可证明
\[
\varphi\circ(k\psi)=k(\varphi\circ\psi). \quad \square
\]

在 $\mathcal{L}(V)$ 中，定义线性变换 $\varphi$ 的 $n$ 次幂为 $n$ 个 $\varphi$ 的复合，则不难验证：
\begin{equation}
\varphi^n\circ\varphi^m=\varphi^{n+m},\quad (\varphi^n)^m=\varphi^{mn}.
\tag{4.2.1}
\end{equation}
若 $\varphi$ 是双射，即为 $V$ 上的自同构，则 $\varphi^{-1}$ 也是 $V$ 上的线性变换 (也是自同构)，称 $\varphi^{-1}$ 为 $\varphi$ 的逆变换. 如定义
\[
\varphi^{-n}=(\varphi^{-1})^n,
\]
则不难验证：
\begin{equation}
\varphi^{-n}=(\varphi^n)^{-1}.
\tag{4.2.2}
\end{equation}
这时定义
\[
\varphi^0=I_V,
\]
则 (4.2.1) 式对一切整数均成立. 但需注意 $\varphi$ 的负数次幂仅对自同构 (又称可逆变换或非异变换) 有意义.

读者需要特别注意的是，线性变换的复合通常不满足交换律，即一般来说，
\[
\varphi\circ\psi\neq\psi\circ\varphi.
\]
因此一般来说，$(\varphi\circ\psi)^n\neq\varphi^n\circ\psi^n$.

如果 $\varphi$ 与 $\psi$ 都是可逆线性变换，则 $\varphi\circ\psi$ 也是可逆线性变换，且
\[
(\varphi\circ\psi)^{-1}=\psi^{-1}\circ\varphi^{-1}.
\]
对任一非零数 $k$，若 $\varphi$ 可逆，则 $k\varphi$ 也可逆，且
\[
(k\varphi)^{-1}=k^{-1}\varphi^{-1}.
\]
读者不难自己验证上述结论.

\subsection*{习题 4.2}

1. 求证：数域 $\mathbb{K}$ 上的 $n$ 阶矩阵全体组成的线性空间在矩阵乘法下是 $\mathbb{K}$ 上的代数.

2. 求证：对 $\mathcal{L}(V)$ 中的任一元 $\varphi$ 及 $\mathbb{K}$ 中的数 $k$，有
\[
k\varphi=(kI_V)\circ\varphi=\varphi\circ(kI_V).
\]

3. 设 $V$ 是实系数多项式全体构成的实线性空间，定义 $V$ 上的变换 $D,S$ 如下：
\[
D(f(x))=\frac{\mathrm{d}}{\mathrm{d}x}f(x),\quad S(f(x))=\int_0^x f(t)\mathrm{d}t.
\]
证明：$D,S$ 均为 $V$ 上的线性变换且 $DS=I_V$，但 $SD\neq I_V$.

4. 若线性变换 $\varphi$ 适合 $\varphi^2=\varphi$，则称 $\varphi$ 是幂等变换. 求证：若 $A^2=A$ (即 $A$ 是幂等阵)，则在上一节例 4.1.3 中定义的线性变换是幂等变换.

5. 设 $\varphi$ 是 $n$ 维线性空间 $V$ 上的线性变换，$\alpha\in V$. 若 $\varphi^{m-1}(\alpha)\neq 0$，而 $\varphi^m(\alpha)=0$，求证：$\alpha,\varphi(\alpha),\varphi^2(\alpha),\cdots,\varphi^{m-1}(\alpha)$ 线性无关.

6. 设 $\varphi$ 是数域 $\mathbb{F}$ 上线性空间 $V$ 上的线性变换，若存在正整数 $n$ 以及 $a_1,a_2,\cdots,a_n\in\mathbb{F}$，使
\[
\varphi^n+a_1\varphi^{n-1}+\cdots+a_{n-1}\varphi+a_nI_V=0,
\]
其中 $I_V$ 表示恒等变换并且 $a_n\neq 0$，求证：$\varphi$ 是 $V$ 上的自同构.

7. 设 $\varphi$ 是 $n$ 维线性空间 $V$ 上的线性变换，证明：$\varphi$ 是可逆变换的充分必要条件是 $\varphi$ 将 $V$ 的基变为基.

\section*{§4.3\quad 线性映射与矩阵}

我们已经定义了线性空间之间的线性映射及其运算. 线性映射是一个比较抽象的“几何”概念，不便于计算和研究. 我们这一节的目的是要将这个抽象的概念“代数化”. 我们在第三章中通过引进线性空间的基，将抽象的线性空间和行向量空间或列向量空间联系了起来. 在例 4.1.3 中，我们也注意到，可以用矩阵来定义列向量空间的线性映射. 因此自然地，我们希望将线性映射和矩阵联系起来.

首先注意到如果取定线性空间 $V$ 的一组基，则从 $V$ 到另一个线性空间 $U$ 的线性映射 $\varphi$ 完全被它在基上的作用所决定，即我们有下面的引理.

\noindent\textbf{引理 4.3.1}\quad 设有 $\mathbb{K}$ 上的线性空间 $V$ 和 $U$，$\{e_1,e_2,\cdots,e_n\}$ 是 $V$ 的一组基.

(1) 如有 $V$ 到 $U$ 的线性映射 $\varphi$ 和 $\psi$，满足 $\psi(e_i)=\varphi(e_i)(i=1,2,\cdots,n)$，则 $\psi=\varphi$;

(2) 给定 $U$ 中 $n$ 个向量 $\beta_1,\beta_2,\cdots,\beta_n$，有且只有一个从 $V$ 到 $U$ 的线性映射 $\varphi$，满足 $\varphi(e_i)=\beta_i(i=1,2,\cdots,n)$.

\noindent\textbf{证明}\quad (1) 对任意的 $\alpha\in V$，设 $\alpha=\lambda_1e_1+\lambda_2e_2+\cdots+\lambda_ne_n$，则
\[
\begin{aligned}
\psi(\alpha)
&=\psi(\lambda_1e_1+\lambda_2e_2+\cdots+\lambda_ne_n)\\
&=\lambda_1\psi(e_1)+\lambda_2\psi(e_2)+\cdots+\lambda_n\psi(e_n)\\
&=\lambda_1\varphi(e_1)+\lambda_2\varphi(e_2)+\cdots+\lambda_n\varphi(e_n)\\
&=\varphi(\lambda_1e_1+\lambda_2e_2+\cdots+\lambda_ne_n)=\varphi(\alpha),
\end{aligned}
\]
因此 $\psi=\varphi$.

(2) 定义 $V\to U$ 的映射 $\varphi$ 如下：对 $V$ 中任意的 $\alpha=\lambda_1e_1+\lambda_2e_2+\cdots+\lambda_ne_n$，
\[
\varphi(\alpha)=\lambda_1\beta_1+\lambda_2\beta_2+\cdots+\lambda_n\beta_n.
\]
容易验证这是 $V$ 到 $U$ 的线性映射，且 $\varphi(e_i)=\beta_i(i=1,2,\cdots,n)$. 唯一性由 (1) 即得. $\square$

现考虑这样一个问题：设 $V$ 与 $U$ 分别是数域 $\mathbb{K}$ 上 $n$ 维及 $m$ 维线性空间，$\{e_1,e_2,\cdots,e_n\}$ 是 $V$ 的基，$\{f_1,f_2,\cdots,f_m\}$ 是 $U$ 的基. 又设 $\varphi$ 是 $V\to U$ 的线性映射且已知 $V$ 的基向量在 $\varphi$ 下的像. 若 $V$ 中向量 $\alpha$ 在给定基下的坐标向量是 $(\lambda_1,\lambda_2,\cdots,\lambda_n)'$，问：如何来求 $\varphi(\alpha)$ 在 $U$ 的基 $\{f_1,f_2,\cdots,f_m\}$ 下的坐标向量？

设
\begin{equation}
\left\{
\begin{array}{l}
\varphi(e_1)=a_{11}f_1+a_{12}f_2+\cdots+a_{1m}f_m,\\
\varphi(e_2)=a_{21}f_1+a_{22}f_2+\cdots+a_{2m}f_m,\\
\quad\cdots\cdots\cdots\\
\varphi(e_n)=a_{n1}f_1+a_{n2}f_2+\cdots+a_{nm}f_m.
\end{array}\right.
\tag{4.3.1}
\end{equation}
因为 $\alpha=\lambda_1e_1+\lambda_2e_2+\cdots+\lambda_ne_n$，故
\[
\begin{aligned}
\varphi(\alpha)
&=\lambda_1\varphi(e_1)+\lambda_2\varphi(e_2)+\cdots+\lambda_n\varphi(e_n)\\
&=\lambda_1\left(\sum_{j=1}^m a_{1j}f_j\right)+\lambda_2\left(\sum_{j=1}^m a_{2j}f_j\right)+\cdots+\lambda_n\left(\sum_{j=1}^m a_{nj}f_j\right)\\
&=\left(\sum_{i=1}^n\lambda_i a_{i1}\right)f_1+\left(\sum_{i=1}^n\lambda_i a_{i2}\right)f_2+\cdots+\left(\sum_{i=1}^n\lambda_i a_{im}\right)f_m.
\end{aligned}
\]

这表明 $\varphi(\alpha)$ 在 $\{f_1,f_2,\cdots,f_m\}$ 下的坐标向量 $(\mu_1,\mu_2,\cdots,\mu_m)'$ 为
\begin{equation}
\left(\begin{array}{c}
\mu_1\\
\mu_2\\
\vdots\\
\mu_m
\end{array}\right)
=
\left(\begin{array}{cccc}
a_{11} & a_{21} & \cdots & a_{n1}\\
a_{12} & a_{22} & \cdots & a_{n2}\\
\vdots & \vdots & & \vdots\\
a_{1m} & a_{2m} & \cdots & a_{nm}
\end{array}\right)
\left(\begin{array}{c}
\lambda_1\\
\lambda_2\\
\vdots\\
\lambda_n
\end{array}\right).
\tag{4.3.2}
\end{equation}
上式中的矩阵 $A=(a_{ji})_{m\times n}$ 是 (4.3.1) 式中系数矩阵的转置. 我们称这个矩阵为 $\varphi$ 在给定基 $\{e_1,e_2,\cdots,e_n\}$ 与 $\{f_1,f_2,\cdots,f_m\}$ 下的表示矩阵，或简称为 $\varphi$ 在给定基下的矩阵.

注意 (4.3.1) 式和 (4.3.2) 式是等价的. 我们在上面从 (4.3.1) 式推出了 (4.3.2) 式. 反过来，如果 (4.3.2) 式成立，即 $V$ 中向量 $\alpha$ 在线性映射 $\varphi$ 的作用下其坐标向量可以用 (4.3.2) 式来表示，则由 $e_1$ 的坐标向量是 $(1,0,\cdots,0)'$ 以及 (4.3.2) 式得到 $\varphi(e_1)$ 的坐标向量是 $(a_{11},a_{12},\cdots,a_{1m})'$ (即表示矩阵的第一个列向量). 因此
\[
\varphi(e_1)=a_{11}f_1+a_{12}f_2+\cdots+a_{1m}f_m.
\]
同理，
\[
\varphi(e_i)=a_{i1}f_1+a_{i2}f_2+\cdots+a_{im}f_m.
\]
于是我们得到了 (4.3.1) 式.

我们在给定基后，由从 $V$ 到 $U$ 的一个线性映射得到了一个 $m\times n$ 矩阵，反过来，给定一个 $\mathbb{K}$ 上的 $m\times n$ 矩阵 $A=(a_{ji})$，由引理 4.3.1，我们可以得到 $V\to U$ 的唯一一个线性映射 $\varphi$，使 $\varphi(e_i)$ 适合 (4.3.1) 式.

若记 $\mathcal{L}(V,U)$ 为从 $V$ 到 $U$ 的线性映射全体组成的集合，$M_{m\times n}(\mathbb{K})$ 是 $\mathbb{K}$ 上 $m\times n$ 矩阵全体组成的集合，则我们得到了一个从 $\mathcal{L}(V,U)$ 到 $M_{m\times n}(\mathbb{K})$ 的映射 $T$，对任意的 $\varphi\in\mathcal{L}(V,U)$，$T(\varphi)=A$，其中 $A$ 是 $\varphi$ 在给定基下的表示矩阵. 前面的分析告诉我们：$T$ 是一个一一对应.

我们已经知道，$M_{m\times n}(\mathbb{K})$ 在矩阵的加法与数乘下是 $\mathbb{K}$ 上的线性空间，$\mathcal{L}(V,U)$ 也是 $\mathbb{K}$ 上的线性空间. 这两个线性空间同构吗？即一一对应 $T$ 保持加法运算和数乘运算吗？

先做一些符号上的说明. 假设 $V$ 的基为 $\{e_1,e_2,\cdots,e_n\}$，$U$ 的基为 $\{f_1,f_2,\cdots,f_m\}$. 记 $\eta_1$ 是 $V$ 到 $\mathbb{K}_n$ 的线性同构：若 $\alpha=a_1e_1+a_2e_2+\cdots+a_ne_n$，则
\[
\eta_1(\alpha)=(a_1,a_2,\cdots,a_n)'.
\]
同样，若 $\beta=b_1f_1+b_2f_2+\cdots+b_mf_m$，则令
\[
\eta_2(\beta)=(b_1,b_2,\cdots,b_m)'.
\]

设 $\varphi\in\mathcal{L}(V,U)$，$T(\varphi)=A$ 是 $\varphi$ 在给定基下的表示矩阵. 我们约定用 $\varphi_A$ 表示在例 4.1.3 中定义的从 $\mathbb{K}_n\to\mathbb{K}_m$ 的线性映射，即若 $x\in\mathbb{K}_n$，则 $\varphi_A(x)=Ax$.

\noindent\textbf{定理 4.3.1}\quad 设 $T$ 是由上面定义的从 $\mathcal{L}(V,U)$ 到 $M_{m\times n}(\mathbb{K})$ 的映射，则 $T$ 是一个线性同构. 不仅如此，$\eta_2\varphi=\varphi_A\eta_1$，即有图 4.1 所示的交换图.

\noindent\textbf{证明}\quad 我们先验证 $T$ 是一个线性映射. 设 $\varphi,\psi$ 是 $V\to U$ 的线性映射且 $T(\varphi)=A=(a_{ji})$，$T(\psi)=B=(b_{ji})$. 对任意的 $e_i(i=1,2,\cdots,n)$，有
\[
\begin{aligned}
(\varphi+\psi)(e_i)
&=\varphi(e_i)+\psi(e_i)\\
&=\sum_{j=1}^m a_{ij}f_j+\sum_{j=1}^m b_{ij}f_j\\
&=\sum_{j=1}^m(a_{ij}+b_{ij})f_j,
\end{aligned}
\]
因此
\[
T(\varphi+\psi)=A+B=T(\varphi)+T(\psi).
\]
同理，可证明对任意的 $k\in\mathbb{K}$ 及 $\varphi\in\mathcal{L}(V,U)$，有
\[
T(k\varphi)=kA=kT(\varphi).
\]
这表明 $T$ 是线性映射. 因为 $T$ 是一一对应，故 $T$ 是线性同构.

要证明图 4.1 所示的交换性，只要对 $V$ 的任一基向量 $e_i$，验证 $\eta_2\varphi(e_i)=\varphi_A\eta_1(e_i)$ 即可. 设
\[
A=
\left(\begin{array}{cccc}
a_{11} & a_{21} & \cdots & a_{n1}\\
a_{12} & a_{22} & \cdots & a_{n2}\\
\vdots & \vdots & & \vdots\\
a_{1m} & a_{2m} & \cdots & a_{nm}
\end{array}\right),
\]
则 $\varphi(e_i)=a_{i1}f_1+a_{i2}f_2+\cdots+a_{im}f_m$，故
\[
\eta_2\varphi(e_i)=
\left(\begin{array}{c}
a_{i1}\\
a_{i2}\\
\vdots\\
a_{im}
\end{array}\right).
\]
注意到 $\eta_1(e_i)$ 是第 $i$ 个标准单位列向量，因此 $\varphi_A\eta_1(e_i)$ 就等于 $A$ 的第 $i$ 个列向量，即 $\varphi_A\eta_1(e_i)=\eta_2\varphi(e_i)$. 图的交换性成立. $\square$

下面的定理告诉我们，矩阵乘法的几何意义是线性映射的复合.

\noindent\textbf{定理 4.3.2}\quad 同定理 4.3.1 的假设，再设 $W$ 是 $\mathbb{K}$ 上的线性空间，$\{g_1,g_2,\cdots,g_p\}$ 是 $W$ 的一组基，$\psi\in\mathcal{L}(U,W)$，则 $T(\psi\varphi)=T(\psi)T(\varphi)$.

\noindent\textbf{证明}\quad 设 $T(\varphi)=A=(a_{ji})_{m\times n}$，$T(\psi)=B=(b_{kj})_{p\times m}$ 分别是 $\varphi,\psi$ 在给定基下的表示矩阵，又 $\alpha=\lambda_1e_1+\lambda_2e_2+\cdots+\lambda_ne_n$ 是 $V$ 中任一向量，则 $\varphi(\alpha)$ 的坐标向量为
\[
\left(\begin{array}{c}
\mu_1\\
\mu_2\\
\vdots\\
\mu_m
\end{array}\right)
=A
\left(\begin{array}{c}
\lambda_1\\
\lambda_2\\
\vdots\\
\lambda_n
\end{array}\right).
\]
$\psi(\varphi(\alpha))$ 的坐标向量为
\[
\left(\begin{array}{c}
\xi_1\\
\xi_2\\
\vdots\\
\xi_p
\end{array}\right)
=B
\left(\begin{array}{c}
\mu_1\\
\mu_2\\
\vdots\\
\mu_m
\end{array}\right).
\]
因此
\[
\left(\begin{array}{c}
\xi_1\\
\xi_2\\
\vdots\\
\xi_p
\end{array}\right)
=BA
\left(\begin{array}{c}
\lambda_1\\
\lambda_2\\
\vdots\\
\lambda_n
\end{array}\right).
\]
这表明 $T(\psi\varphi)=BA=T(\psi)T(\varphi)$. $\square$

现在我们考察 $V$ 上全体线性变换 $\mathcal{L}(V)$. 取定 $V$ 的一组基 $\{e_1,e_2,\cdots,e_n\}$，对任一 $\varphi\in\mathcal{L}(V)$，设
\begin{equation}
\left\{
\begin{array}{l}
\varphi(e_1)=a_{11}e_1+a_{12}e_2+\cdots+a_{1n}e_n,\\
\varphi(e_2)=a_{21}e_1+a_{22}e_2+\cdots+a_{2n}e_n,\\
\quad\cdots\cdots\cdots\\
\varphi(e_n)=a_{n1}e_1+a_{n2}e_2+\cdots+a_{nn}e_n,
\end{array}\right.
\tag{4.3.3}
\end{equation}
则 $\varphi$ 在基 $\{e_1,e_2,\cdots,e_n\}$ 下的表示矩阵定义为 $n$ 阶方阵 $A=(a_{ji})$，即 (4.3.3) 式中系数矩阵的转置，记为 $T(\varphi)=A$.

\noindent\textbf{定理 4.3.3}\quad $T:\mathcal{L}(V)\to M_n(\mathbb{K})$ 是线性同构，并对任意的 $\varphi,\psi\in\mathcal{L}(V)$，有
\[
T(\psi\varphi)=T(\psi)T(\varphi),
\]
即 $T$ 保持了乘法.

\noindent\textbf{证明}\quad 由类似于定理 4.3.1 和定理 4.3.2 的证明可得结论. $\square$

\noindent\textbf{推论 4.3.1}\quad 上述同构 $T$ 有下列性质：

(1) $T(I_V)=I_n$;

(2) $\varphi$ 是 $V$ 上自同构的充分必要条件是 $T(\varphi)$ 为可逆阵且这时有
\[
T(\varphi^{-1})=T(\varphi)^{-1}.
\]

\noindent\textbf{证明}\quad (1) 当 $\varphi=I_V$ 时，(4.3.3) 式中的系数矩阵为 $I_n$，其转置也为 $I_n$，因此结论成立.

(2) 由 $\varphi\varphi^{-1}=I_V$，得
\[
T(\varphi)T(\varphi^{-1})=T(\varphi\varphi^{-1})=T(I_V)=I_n.
\]
此即 $T(\varphi^{-1})=T(\varphi)^{-1}$. 充分性也不难验证. $\square$

上面的这些结论在线性映射与矩阵之间建立起了桥梁，它使我们能用代数的工具 (矩阵) 来研究几何的对象 (线性映射)，也能用几何的方法 (线性映射) 来研究代数的对象 (矩阵).

我们继续讨论线性变换的问题. 我们知道线性变换的表示矩阵是与线性空间中的基联系在一起的. 一般来说，当基发生变化时，同一个线性变换在不同基下的表示矩阵是不相同的. 如果我们已经知道了两组基及其过渡矩阵，同一个线性变换在这两组基下的表示矩阵有什么关系呢？

\noindent\textbf{定理 4.3.4}\quad 设 $V$ 是数域 $\mathbb{K}$ 上的线性空间，$\varphi\in\mathcal{L}(V)$，又设 $\{e_1,e_2,\cdots,e_n\}$ 及 $\{f_1,f_2,\cdots,f_n\}$ 是 $V$ 的两组基且从 $\{e_1,e_2,\cdots,e_n\}$ 到 $\{f_1,f_2,\cdots,f_n\}$ 的过渡矩阵为 $P$. 若 $\varphi$ 在基 $\{e_1,e_2,\cdots,e_n\}$ 下的表示矩阵为 $A$，在基 $\{f_1,f_2,\cdots,f_n\}$ 下的表示矩阵为 $B$，则
\[
B=P^{-1}AP.
\]

\noindent\textbf{证明}\quad 设 $\alpha$ 是 $V$ 中任一向量且
\[
\alpha=\lambda_1e_1+\lambda_2e_2+\cdots+\lambda_ne_n=\mu_1f_1+\mu_2f_2+\cdots+\mu_nf_n,
\]
则由第三章中的 (3.8.3) 式得
\begin{equation}
\left(\begin{array}{c}
\lambda_1\\
\lambda_2\\
\vdots\\
\lambda_n
\end{array}\right)
=P
\left(\begin{array}{c}
\mu_1\\
\mu_2\\
\vdots\\
\mu_n
\end{array}\right).
\tag{4.3.4}
\end{equation}
设
\begin{equation}
\varphi(\alpha)=\xi_1e_1+\xi_2e_2+\cdots+\xi_ne_n=\eta_1f_1+\eta_2f_2+\cdots+\eta_nf_n,
\tag{4.3.5}
\end{equation}
则由 (4.3.2) 式得
\begin{equation}
\left(\begin{array}{c}
\xi_1\\
\xi_2\\
\vdots\\
\xi_n
\end{array}\right)
=A
\left(\begin{array}{c}
\lambda_1\\
\lambda_2\\
\vdots\\
\lambda_n
\end{array}\right),\quad
\left(\begin{array}{c}
\eta_1\\
\eta_2\\
\vdots\\
\eta_n
\end{array}\right)
=B
\left(\begin{array}{c}
\mu_1\\
\mu_2\\
\vdots\\
\mu_n
\end{array}\right).
\tag{4.3.6}
\end{equation}
另一方面由 (4.3.5) 式及第三章的 (3.8.3) 式有
\begin{equation}
\left(\begin{array}{c}
\xi_1\\
\xi_2\\
\vdots\\
\xi_n
\end{array}\right)
=P
\left(\begin{array}{c}
\eta_1\\
\eta_2\\
\vdots\\
\eta_n
\end{array}\right).
\tag{4.3.7}
\end{equation}
由 (4.3.4) 式、(4.3.6) 式和 (4.3.7) 式得
\begin{equation}
PB
\left(\begin{array}{c}
\mu_1\\
\mu_2\\
\vdots\\
\mu_n
\end{array}\right)
=AP
\left(\begin{array}{c}
\mu_1\\
\mu_2\\
\vdots\\
\mu_n
\end{array}\right).
\tag{4.3.8}
\end{equation}

但 $\alpha$ 是任意的，即 (4.3.8) 式中的 $\mu_i$ 是任意的，因此
\[
PB=AP,
\]
即
\[
B=P^{-1}AP. \quad \square
\]

\noindent\textbf{定义 4.3.1}\quad 若 $A,B$ 为 $n$ 阶方阵且存在 $n$ 阶非异阵 $P$，使
\[
B=P^{-1}AP,
\]
则称 $A$ 与 $B$ 相似，记为 $A\approx B$.

定理 4.3.4 表明：$V$ 上的线性变换 $\varphi$ 在不同基下的表示矩阵是相似的.

\noindent\textbf{命题 4.3.1}\quad 相似关系是一种等价关系，即

(1) $A\approx A$;

(2) 若 $A\approx B$，则 $B\approx A$;

(3) 若 $A\approx B,B\approx C$，则 $A\approx C$.

\noindent\textbf{证明}\quad (1) $A=I_n^{-1}AI_n$，故 $A\approx A$.

(2) 若 $B=P^{-1}AP$，则 $A=PBP^{-1}$，因此 $B\approx A$.

(3) 若 $B=P^{-1}AP$，$C=Q^{-1}BQ$，则 $C=(PQ)^{-1}A(PQ)$，故 $A\approx C$. $\square$

定理 4.3.4 揭示了同一线性变换在不同基下表示矩阵之间的关系. 一个十分重要的问题是：对一个线性变换 $\varphi$ 能否找到一组适当的基，使 $\varphi$ 在这组基下的表示矩阵具有比较简单的形状？这个问题的代数提法是：给定一个 $n$ 阶矩阵 $A$，能否找到一种方法，使得 $A$ 相似于一个比较简单的矩阵？我们将在第六章和第七章探讨这一问题.

\subsection*{习题 4.3}

1. 设 $V$ 是实数域上次数小于 $n$ 的一元多项式全体组成的线性空间，$\varphi$ 为多项式的求导运算. 试求 $\varphi$ 在 $V$ 的基 $\{1,x,\cdots,x^{n-1}\}$ 下的表示矩阵.

2. 设 $\mathbb{K}^4$ 上的线性变换 $\varphi$ 在一组基 $\{e_1,e_2,e_3,e_4\}$ 下的表示矩阵为
\[
\left(\begin{array}{rrrr}
1 & 2 & 3 & 2\\
-1 & 0 & 3 & 1\\
2 & 1 & 5 & -1\\
1 & 1 & 2 & 2
\end{array}\right),
\]
试求 $\varphi$ 在下列基下的表示矩阵：

(1) $\{e_4,e_3,e_2,e_1\}$;

(2) $\{e_1,e_1+e_2,e_1+e_2+e_3,e_1+e_2+e_3+e_4\}$.

3. 设 $V$ 是 Descartes 平面，试求绕原点逆时针旋转 $\theta$ 角的线性变换在基 $\{(1,0),(0,1)\}$ 下的表示矩阵.

4. 设 $\varphi,\psi$ 是线性空间 $V$ 上的线性变换，且 $\varphi$ 是可逆变换，设 $\varphi,\psi$ 在 $V$ 的第一组基下的表示矩阵分别为 $A,B$，又第一组基到第二组基的过渡矩阵为 $P$，试求线性变换 $\psi\varphi^{-2}+2\varphi+I_V$ 在 $V$ 的第二组基下的表示矩阵.

5. 设 $V,U$ 分别是数域 $\mathbb{K}$ 上的 $n,m$ 维线性空间，求证：$\mathcal{L}(V,U)$ 是 $mn$ 维线性空间.

6. 证明：相似的矩阵具有相同的迹和行列式，即迹和行列式是矩阵在相似关系下的不变量.

7. 设 $A,B$ 是 $n$ 阶方阵且 $A$ 可逆，求证：$AB$ 与 $BA$ 相似.

8. 试求出只与自己相似的所有 $n$ 阶矩阵.

9. 若 $A$ 与 $B$ 相似，$C$ 与 $D$ 相似，证明下列分块矩阵也相似：
\[
\left(\begin{array}{cc}
A & O\\
O & C
\end{array}\right),\quad
\left(\begin{array}{cc}
B & O\\
O & D
\end{array}\right).
\]

10. 设 $A$ 为数域 $\mathbb{K}$ 上的 $n$ 阶方阵，求证：以下三种变换都是相似变换，称为相似初等变换：

(1) 对换 $A$ 的第 $i$ 行与第 $j$ 行，再对换第 $i$ 列与第 $j$ 列；

(2) 将 $A$ 的第 $i$ 行乘以非零常数 $c$，再将第 $i$ 列乘以 $c^{-1}$；

(3) 将 $A$ 的第 $i$ 行乘以常数 $c$ 加到第 $j$ 行上，再将第 $j$ 列乘以 $-c$ 加到第 $i$ 列上.

11. 证明：任一相似变换都是若干次相似初等变换的复合.

12. 设 $A$ 是具有相同行列分块方式的分块矩阵，求证：以下三种变换都是相似变换，称为相似分块初等变换：

(1) 对换 $A$ 的第 $i$ 分块行与第 $j$ 分块行，再对换第 $i$ 分块列与第 $j$ 分块列；

(2) 将 $A$ 的第 $i$ 分块行左乘非异阵 $M$，再将第 $i$ 分块列右乘 $M^{-1}$；

(3) 将 $A$ 的第 $i$ 分块行左乘矩阵 $M$ 加到第 $j$ 分块行上，再将第 $j$ 分块列右乘 $-M$ 加到第 $i$ 分块列上.

13. 证明定理 4.3.4 的逆命题：若 $n$ 阶方阵 $A$ 与 $B$ 相似，则它们可看成是同一个线性变换在两组基下的表示矩阵.

14. 设 $\varphi$ 是线性空间 $V$ 到 $U$ 的线性映射，$\{e_1,e_2,\cdots,e_n\}$ 和 $\{f_1,f_2,\cdots,f_n\}$ 是 $V$ 的两组基，$\{e_1,e_2,\cdots,e_n\}$ 到 $\{f_1,f_2,\cdots,f_n\}$ 的过渡矩阵为 $P$. $\{g_1,g_2,\cdots,g_m\}$ 和 $\{h_1,h_2,\cdots,h_m\}$ 是 $U$ 的两组基，$\{g_1,g_2,\cdots,g_m\}$ 到 $\{h_1,h_2,\cdots,h_m\}$ 的过渡矩阵为 $Q$. 又设 $\varphi$ 在基 $\{e_1,e_2,\cdots,e_n\}$ 和基 $\{g_1,g_2,\cdots,g_m\}$ 下的表示矩阵为 $A$，在基 $\{f_1,f_2,\cdots,f_n\}$ 和基 $\{h_1,h_2,\cdots,h_m\}$ 下的表示矩阵为 $B$. 求证：$B=Q^{-1}AP$.

\section*{§4.4\quad 线性映射的像与核}

\noindent\textbf{定义 4.4.1}\quad 设 $\varphi$ 是数域 $\mathbb{K}$ 上线性空间 $V$ 到 $U$ 的线性映射，$\varphi$ 的全体像元素组成 $U$ 的子集称为 $\varphi$ 的像，记为 $\operatorname{Im}\varphi$. 又，$V$ 中在 $\varphi$ 下映射为零向量的全体向量构成 $V$ 的子集，称为 $\varphi$ 的核，记为 $\operatorname{Ker}\varphi$.

\noindent\textbf{命题 4.4.1}\quad 设 $\varphi$ 是线性空间 $V\to U$ 的线性映射，则 $\operatorname{Im}\varphi$ 是 $U$ 的子空间，$\operatorname{Ker}\varphi$ 是 $V$ 的子空间.

\noindent\textbf{证明}\quad 设 $\alpha,\beta\in\operatorname{Im}\varphi$，则有 $V$ 中向量 $u,v$，使
\[
\alpha=\varphi(u),\quad \beta=\varphi(v).
\]
由于 $\varphi(u+v)=\varphi(u)+\varphi(v)=\alpha+\beta$，故 $\alpha+\beta\in\operatorname{Im}\varphi$. 若 $k\in\mathbb{K}$，则 $\varphi(ku)=k\varphi(u)=k\alpha$，故 $k\alpha\in\operatorname{Im}\varphi$. 因此 $\operatorname{Im}\varphi$ 是 $U$ 的子空间.

设 $u,v\in\operatorname{Ker}\varphi$，则 $\varphi(u)=\varphi(v)=0$，从而
\[
\varphi(u+v)=\varphi(u)+\varphi(v)=0,
\]
这说明 $u+v\in\operatorname{Ker}\varphi$. 类似地可证明 $ku\in\operatorname{Ker}\varphi$. 因此 $\operatorname{Ker}\varphi$ 是 $V$ 的子空间. $\square$

\noindent\textbf{推论 4.4.1}\quad 线性映射 $\varphi$ 是满映射的充分必要条件是 $\dim\operatorname{Im}\varphi=\dim U$；线性映射 $\varphi$ 是单映射的充分必要条件是 $\operatorname{Ker}\varphi=0$.

\noindent\textbf{证明}\quad 第一个结论显然成立. 对第二个结论，设 $\varphi$ 是单映射，若 $v\in\operatorname{Ker}\varphi$，则 $\varphi(v)=0=\varphi(0)$，于是 $v=0$，即 $\operatorname{Ker}\varphi=0$. 反之，设 $\operatorname{Ker}\varphi=0$，若 $\varphi(u)=\varphi(v)$，则 $\varphi(u-v)=\varphi(u)-\varphi(v)=0$，故 $u-v=0$，即 $u=v$. $\square$

\noindent\textbf{定义 4.4.2}\quad 设 $\varphi$ 是 $V\to U$ 的线性映射，像空间 $\operatorname{Im}\varphi$ 的维数称为 $\varphi$ 的秩，记作 $\operatorname{r}(\varphi)$，核空间 $\operatorname{Ker}\varphi$ 的维数称为 $\varphi$ 的零度.

如果已知一个线性映射的表示矩阵，那么它的像空间和核空间的维数如何确定？像空间和核空间如何用已知的基向量来表示？在回答这些问题之前，我们先通过一个引理引入线性映射的限制的概念，这一概念在后面将会经常用到.

\noindent\textbf{引理 4.4.1}\quad 设 $\varphi:V\to U$ 为线性映射，$V'\subseteq V$，$U'\subseteq U$ 为子空间且满足 $\varphi(V')\subseteq U'$，则通过定义域的限制可得线性映射 $\varphi':V'\to U'$，使得 $\varphi'$ 与 $\varphi$ 具有相同的映射法则. 进一步，若 $\varphi$ 是单映射，则 $\varphi'$ 也是单映射.

\noindent\textbf{证明}\quad 定义 $\varphi':V'\to U'$ 如下：对任一 $v'\in V'$，$\varphi'(v')=\varphi(v')\in U'$. 显然 $\varphi'$ 是一个定义好的映射，它其实是将 $\varphi$ 的定义域限制在 $V'$ 上得到的映射，当然与 $\varphi$ 具有相同的映射法则. 由 $\varphi$ 是线性映射容易验证 $\varphi'$ 也是线性映射. 注意到 $\operatorname{Ker}\varphi'=\operatorname{Ker}\varphi\cap V'$，因此第二个结论由推论 4.4.1 即得. $\square$

\noindent\textbf{例 4.4.1}\quad 设 $V$ 是 Descartes 平面，$\varphi$ 是绕原点逆时针旋转 $\theta$ 角的线性变换. 设 $V'$ 是 $x$-轴所在的一维子空间，$U'$ 是 $\theta$ 角直线所在的一维子空间，则限制映射 $\varphi':V'\to U'$ 不仅是单线性映射，也是满线性映射.

\noindent\textbf{定理 4.4.1}\quad 设 $V,U$ 分别是数域 $\mathbb{K}$ 上的 $n$ 维和 $m$ 维线性空间，又设 $\{e_1,e_2,\cdots,e_n\}$ 是 $V$ 的基，$\{f_1,f_2,\cdots,f_m\}$ 是 $U$ 的基. 设 $\varphi$ 是 $V\to U$ 的线性映射，它在给定基下的表示矩阵为 $A$，则
\[
\dim\operatorname{Im}\varphi=\operatorname{rank}(A),\quad \dim\operatorname{Ker}\varphi=n-\operatorname{rank}(A).
\]

\noindent\textbf{证明}\quad 我们沿用定理 4.3.1 中的记号和术语. 考虑定理 4.3.1 中的交换图 4.1. 我们先证明 $\eta_1(\operatorname{Ker}\varphi)\subseteq\operatorname{Ker}\varphi_A$，$\eta_2(\operatorname{Im}\varphi)\subseteq\operatorname{Im}\varphi_A$. 任取 $v\in\operatorname{Ker}\varphi$，则由图 4.1 的交换性可得 $\varphi_A(\eta_1(v))=\eta_2(\varphi(v))=\eta_2(0)=0$，即 $\eta_1(v)\in\operatorname{Ker}\varphi_A$，从而 $\eta_1(\operatorname{Ker}\varphi)\subseteq\operatorname{Ker}\varphi_A$. 任取 $u\in\operatorname{Im}\varphi$，则存在 $v\in V$ 使得 $u=\varphi(v)$，由图 4.1 的交换性可得 $\eta_2(u)=\eta_2(\varphi(v))=\varphi_A(\eta_1(v))\in\operatorname{Im}\varphi_A$，从而 $\eta_2(\operatorname{Im}\varphi)\subseteq\operatorname{Im}\varphi_A$.

注意到 $\eta_1,\eta_2$ 都是线性同构，由引理 4.4.1 通过定义域的限制可以得到两个单线性映射 $\eta'_1:\operatorname{Ker}\varphi\to\operatorname{Ker}\varphi_A$ 和 $\eta'_2:\operatorname{Im}\varphi\to\operatorname{Im}\varphi_A$. 接下去我们证明这两个线性映射也是满射. 任取 $\alpha\in\operatorname{Ker}\varphi_A$，因为 $\eta_1$ 是一一对应，故存在 $v\in V$，使得 $\eta_1(v)=\alpha$，于是由图 4.1 的交换性可得 $0=\varphi_A(\alpha)=\varphi_A(\eta_1(v))=\eta_2(\varphi(v))$. 因为 $\eta_2$ 也是一一对应，故 $\varphi(v)=0$，即 $v\in\operatorname{Ker}\varphi$，于是 $\eta'_1(v)=\eta_1(v)=\alpha$，即 $\eta'_1:\operatorname{Ker}\varphi\to\operatorname{Ker}\varphi_A$ 是满射，从而是线性同构. 任取 $\beta\in\operatorname{Im}\varphi_A$，则存在 $\alpha\in\mathbb{K}_n$，使得 $\varphi_A(\alpha)=\beta$. 因为 $\eta_1$ 是一一对应，故存在 $v\in V$，使得 $\eta_1(v)=\alpha$，于是由图 4.1 的交换性可得 $\beta=\varphi_A(\alpha)=\varphi_A(\eta_1(v))=\eta_2(\varphi(v))$. 令 $u=\varphi(v)\in\operatorname{Im}\varphi$，则 $\eta'_2(u)=\eta_2(u)=\beta$，即 $\eta'_2:\operatorname{Im}\varphi\to\operatorname{Im}\varphi_A$ 是满射，从而是线性同构.

将 $A$ 表示成列分块的形式：
\[
A=(\alpha_1,\alpha_2,\cdots,\alpha_n),
\]
其中 $\alpha_j$ 是 $A$ 的第 $j$ 列向量. 任取 $x=(x_1,x_2,\cdots,x_n)'\in\mathbb{K}_n$，则 $\operatorname{Im}\varphi_A$ 中的向量 $\varphi_A(x)=Ax=x_1\alpha_1+x_2\alpha_2+\cdots+x_n\alpha_n$，其中 $x_1,x_2,\cdots,x_n$ 可取 $\mathbb{K}$ 中的任意数，因此 $\operatorname{Im}\varphi_A=L(\alpha_1,\alpha_2,\cdots,\alpha_n)$. 由定理 3.9.1 知 $\dim\operatorname{Im}\varphi_A=\operatorname{rank}(A)$.

因为 $\eta'_2:\operatorname{Im}\varphi\to\operatorname{Im}\varphi_A$ 是线性同构，故 $\dim\operatorname{Im}\varphi=\operatorname{rank}(A)$. 又 $\operatorname{Ker}\varphi_A=\{x\in\mathbb{K}_n\mid \varphi_A(x)=Ax=0\}$，即 $\operatorname{Ker}\varphi_A$ 是齐次线性方程组 $Ax=0$ 的解空间，由定理 3.10.2 知 $\dim\operatorname{Ker}\varphi_A=n-\operatorname{rank}(A)$. 因为 $\eta'_1:\operatorname{Ker}\varphi\to\operatorname{Ker}\varphi_A$ 是线性同构，故 $\dim\operatorname{Ker}\varphi=n-\operatorname{rank}(A)$. $\square$

由定理 4.4.1 即得如下两个重要的推论.

\noindent\textbf{推论 4.4.2 (线性映射维数公式)}\quad 设 $\varphi$ 是 $\mathbb{K}$ 上 $n$ 维线性空间 $V$ 到 $\mathbb{K}$ 上 $m$ 维线性空间 $U$ 的线性映射，则
\[
\dim\operatorname{Ker}\varphi+\dim\operatorname{Im}\varphi=\dim V.
\]

\noindent\textbf{推论 4.4.3}\quad 记号同上，$\varphi$ 是满映射的充分必要条件是 $\operatorname{r}(A)=m$，即表示矩阵 $A$ 是一个行满秩阵；$\varphi$ 是单映射的充分必要条件是 $\operatorname{r}(A)=n$，即 $A$ 是一个列满秩阵.

\noindent\textbf{推论 4.4.4}\quad $n$ 维线性空间 $V$ 上的线性变换 $\varphi$ 是可逆变换的充分必要条件为它是单映射或它是满映射.

\noindent\textbf{证明}\quad 若 $\varphi$ 是单映射，则 $\operatorname{Ker}\varphi=0$. 由线性映射维数公式可得 $\dim\operatorname{Im}\varphi=n$，即 $\varphi$ 是满映射，从而 $\varphi$ 是可逆变换 (即自同构).

若 $\varphi$ 是满映射，则 $\dim\operatorname{Im}\varphi=n$，由线性映射维数公式可得 $\operatorname{Ker}\varphi=0$，即 $\varphi$ 是单映射，从而也是自同构. $\square$

\noindent\textbf{推论 4.4.5}\quad $n$ 维线性空间 $V$ 上的线性变换 $\varphi$ 是单映射 (或满映射) 的充分必要条件为它在 $V$ 的任意一组基下的表示矩阵是可逆阵.

\noindent\textbf{证明}\quad 由推论 4.3.1 和推论 4.4.4 即得. 也可用代数方法来证明. 注意到一个 $n$ 阶方阵 $A$ 可逆的充分必要条件是 $A$ 为行满秩阵或列满秩阵，故由推论 4.4.3 即得结论. $\square$

下面的例子将告诉我们如何计算像空间和核空间.

\noindent\textbf{例 4.4.2}\quad 设 $V$ 是 $\mathbb{K}$ 上五维空间，$\{e_1,e_2,e_3,e_4,e_5\}$ 是 $V$ 的基，$U$ 是 $\mathbb{K}$ 上四维空间，$\{f_1,f_2,f_3,f_4\}$ 是 $U$ 的基，线性映射 $\varphi:V\to U$ 在上述基下的表示矩阵为
\[
A=
\left(\begin{array}{rrrrr}
1 & 2 & 1 & -3 & 2\\
2 & 1 & 1 & 1 & -4\\
1 & 1 & 2 & 2 & -2\\
2 & 3 & -5 & -17 & 8
\end{array}\right),
\]

求 $\operatorname{Im}\varphi$ 和 $\operatorname{Ker}\varphi$.

\noindent\textbf{解}\quad 对矩阵 $A$ 进行初等行变换：
\[
\begin{aligned}
A&\to
\left(\begin{array}{rrrrr}
1&2&1&-3&2\\
0&-3&-1&7&-8\\
0&-1&1&5&-4\\
0&-1&-7&-11&4
\end{array}\right)
\to
\left(\begin{array}{rrrrr}
1&2&1&-3&2\\
0&-1&1&5&-4\\
0&-3&-1&7&-8\\
0&-1&-7&-11&4
\end{array}\right)\\
&\to
\left(\begin{array}{rrrrr}
1&0&3&7&-6\\
0&-1&1&5&-4\\
0&0&-4&-8&4\\
0&0&-8&-16&8
\end{array}\right)
\to
\left(\begin{array}{rrrrr}
1&0&3&7&-6\\
0&1&-1&-5&4\\
0&0&1&2&-1\\
0&0&0&0&0
\end{array}\right)\\
&\to
\left(\begin{array}{rrrrr}
1&0&0&1&-3\\
0&1&0&-3&3\\
0&0&1&2&-1\\
0&0&0&0&0
\end{array}\right).
\end{aligned}
\]
因此 $\operatorname{r}(A)=3$，即 $\dim\operatorname{Im}\varphi=3$. 从上面可以看出 $A$ 的前 3 个列向量线性无关，因此它们可以组成 $\operatorname{Im}\varphi$ 的一组基，故
\[
\operatorname{Im}\varphi=k_1(f_1+2f_2+f_3+2f_4)+k_2(2f_1+f_2+f_3+3f_4)+k_3(f_1+f_2+2f_3-5f_4),
\]
其中 $k_i$ 可取 $\mathbb{K}$ 中的任意数.

方程 $Ax=0$ 的基础解系为
\[
\alpha_1=\left(\begin{array}{r}
-1\\
3\\
-2\\
1\\
0
\end{array}\right),\quad
\alpha_2=\left(\begin{array}{r}
3\\
-3\\
1\\
0\\
1
\end{array}\right).
\]
因此
\[
\operatorname{Ker}\varphi=k_1(-e_1+3e_2-2e_3+e_4)+k_2(3e_1-3e_2+e_3+e_5),
\]
其中 $k_i$ 可取 $\mathbb{K}$ 中的任意数.

\subsection*{习题 4.4}

1. 设线性空间 $V$ 上的线性变换 $\varphi$ 在一组基 $\{e_1,e_2,e_3,e_4\}$ 下的表示矩阵为
\[
A=\left(\begin{array}{rrrr}
1 & 0 & 2 & 1\\
-1 & 2 & 1 & 3\\
1 & 2 & 5 & 5\\
2 & -2 & 1 & -2
\end{array}\right),
\]
求 $\varphi$ 的核空间与像空间 (用基的线性组合来表示).

2. 设 $V=V_1\oplus V_2$，$V$ 上的线性变换 $\varphi$ 定义为：
\[
\varphi(v_1+v_2)=v_1,\quad v_1\in V_1,\ v_2\in V_2.
\]
求证：$\varphi^2=\varphi$，并求 $\operatorname{Im}\varphi$ 及 $\operatorname{Ker}\varphi$. 又若 $\{e_1,\cdots,e_r\}$ 是 $V_1$ 的基，$\{e_{r+1},\cdots,e_n\}$ 是 $V_2$ 的基，求 $\varphi$ 在 $V$ 的基 $\{e_1,\cdots,e_r,e_{r+1},\cdots,e_n\}$ 下的表示矩阵.

3. 设 $V$ 是数域 $\mathbb{K}$ 上二阶矩阵全体组成的向量空间，定义 $V$ 上的线性变换 $\varphi$ 如下：
\[
\varphi(A)=
\left(\begin{array}{cc}
1 & 1\\
1 & 1
\end{array}\right)
A
\left(\begin{array}{cc}
2 & 0\\
0 & 1
\end{array}\right),
\]
求 $\varphi$ 的秩和零度.

4. 设 $V=M_n(\mathbb{K})$，对任一 $A\in V$，定义 $\varphi(A)=\operatorname{tr}(A)$. 求证：$\varphi$ 是 $V\to\mathbb{K}$ 的线性映射，并求 $\operatorname{Ker}\varphi$ 的维数及其一组基.

5. 设 $\varphi$ 是有限维线性空间 $V$ 到 $U$ 的线性映射，且 $V$ 的维数大于 $U$ 的维数，求证：$\operatorname{Ker}\varphi\neq 0$.

6. 设 $\varphi$ 是有限维线性空间 $V\to V'$ 的线性映射，$U$ 是 $V'$ 的子空间且 $U\subseteq\operatorname{Im}\varphi$，求证：
\[
\varphi^{-1}(U)=\{v\in V\mid \varphi(v)\in U\}
\]
是 $V$ 的子空间，且
\[
\dim U+\dim\operatorname{Ker}\varphi=\dim\varphi^{-1}(U).
\]

7. 设 $U$ 是有限维线性空间 $V$ 的子空间，$\varphi$ 是 $V$ 上的线性变换，求证：

(1) $\dim U-\dim\operatorname{Ker}\varphi\leq \dim\varphi(U)\leq \dim U$;

(2) $\dim\varphi^{-1}(U)\leq \dim U+\dim\operatorname{Ker}\varphi$.

8. 利用上题证明：若 $A,B$ 是数域 $\mathbb{F}$ 上两个 $n$ 阶方阵，则
\[
\operatorname{r}(A)+\operatorname{r}(B)-n\leq \operatorname{r}(AB)\leq \min\{\operatorname{r}(A),\operatorname{r}(B)\}.
\]

9. 举例说明推论 4.4.4 对无限维线性空间不成立，即存在无限维线性空间 $V$ 上的线性变换 $\varphi$，使得 $\varphi$ 是单映射或满映射，但不是自同构.

\section*{§4.5\quad 不变子空间}

设 $V$ 是数域 $\mathbb{K}$ 上的线性空间，$\varphi$ 是 $V$ 上的线性变换. 现在我们来研究由 $\varphi$ 决定的一类子空间——不变子空间. 不变子空间的理论对以后我们将要学习的相似标准型理论有重要的意义.

\noindent\textbf{定义 4.5.1}\quad 设 $\varphi$ 是线性空间 $V$ 上的线性变换，$U$ 是 $V$ 的子空间，若 $U$ 适合条件
\[
\varphi(U)\subseteq U,
\]
则称 $U$ 是 $\varphi$ 的不变子空间 (或 $\varphi$-不变子空间). 这时把 $\varphi$ 的定义域限制在 $U$ 上，则 $\varphi$ 在 $U$ 上定义了一个线性变换，称为由 $\varphi$ 诱导出的线性变换，或称为 $\varphi$ 在 $U$ 上的限制，记为 $\varphi|_U$.

线性空间 $V$ 上任一线性变换 $\varphi$ 至少有两个不变子空间：零子空间及全空间 $V$. 因此我们把零子空间及全空间 $V$ 称为平凡的 $\varphi$-不变子空间.

\noindent\textbf{例 4.5.1}\quad 线性变换 $\varphi$ 的像与核都是 $\varphi$ 的不变子空间.

\noindent\textbf{证明}\quad $\varphi(\operatorname{Im}\varphi)\subseteq\varphi(V)=\operatorname{Im}\varphi$，$\varphi(\operatorname{Ker}\varphi)=0\subseteq\operatorname{Ker}\varphi$. 由此即得结论. $\square$

\noindent\textbf{例 4.5.2}\quad Descartes 平面上绕原点的旋转，当旋转角 $\theta\neq k\pi$ ($k$ 为整数) 时，没有一维的不变子空间，因此没有非平凡的不变子空间.

\noindent\textbf{例 4.5.3}\quad 设 $\varphi$ 是 $V$ 上的数乘变换，即存在常数 $k$，使 $\varphi(\alpha)=k\alpha$，则 $V$ 的任一子空间都是 $\varphi$ 的不变子空间.

下面我们要讨论线性变换的不变子空间和该线性变换的表示矩阵之间的关系.

\noindent\textbf{定理 4.5.1}\quad 设 $U$ 是 $V$ 上线性变换 $\varphi$ 的不变子空间，且设 $U$ 的基为 $\{e_1,e_2,\cdots,e_r\}$. 将 $\{e_1,e_2,\cdots,e_r\}$ 扩充为 $V$ 的一组基 $\{e_1,e_2,\cdots,e_r,e_{r+1},\cdots,e_n\}$，则 $\varphi$ 在这组基下的表示矩阵具有下列形状：
\begin{equation}
\left(\begin{array}{cccccc}
a_{11} & \cdots & a_{r1} & a_{r+1,1} & \cdots & a_{n1}\\
\vdots & & \vdots & \vdots & & \vdots\\
a_{1r} & \cdots & a_{rr} & a_{r+1,r} & \cdots & a_{nr}\\
0 & \cdots & 0 & a_{r+1,r+1} & \cdots & a_{n,r+1}\\
\vdots & & \vdots & \vdots & & \vdots\\
0 & \cdots & 0 & a_{r+1,n} & \cdots & a_{nn}
\end{array}\right).
\tag{4.5.1}
\end{equation}

\noindent\textbf{证明}\quad 由于 $\varphi(e_i)\in U(i=1,2,\cdots,r)$，因此
\[
\varphi(e_1)=a_{11}e_1+a_{12}e_2+\cdots+a_{1r}e_r,
\]
\[
\varphi(e_2)=a_{21}e_1+a_{22}e_2+\cdots+a_{2r}e_r,
\]
\[
\cdots\cdots\cdots
\]
\[
\varphi(e_r)=a_{r1}e_1+a_{r2}e_2+\cdots+a_{rr}e_r,
\]
即在 $\varphi(e_1),\cdots,\varphi(e_r)$ 的表示式中，$e_{r+1},\cdots,e_n$ 前的系数均为零，因此 $\varphi$ 的表示矩阵具有所要求的形状. $\square$

上述定理的逆命题也是成立的. 即若 $V$ 有一组基 $\{e_1,e_2,\cdots,e_n\}$，使得线性变换 $\varphi$ 在这组基下的表示矩阵是分块上三角阵 (4.5.1)，则由基向量 $e_1,\cdots,e_r$ 生成的子空间 $V_1$ 是 $\varphi$ 的不变子空间. 事实上，由矩阵 (4.5.1) 即知 $\varphi(e_i)\in V_1(i=1,\cdots,r)$. 因此 $\varphi(V_1)\subseteq V_1$.

\noindent\textbf{推论 4.5.1}\quad 设 $V=V_1\oplus V_2$ 且 $V_1,V_2$ 都是线性变换 $\varphi$ 的不变子空间. 又 $\{e_1,\cdots,e_r\}$ 是 $V_1$ 的基，$\{e_{r+1},\cdots,e_n\}$ 是 $V_2$ 的基，则 $\varphi$ 在基 $\{e_1,e_2,\cdots,e_n\}$ 下的表示矩阵为分块对角阵
\begin{equation}
\left(\begin{array}{cc}
A_1 & O\\
O & A_2
\end{array}\right),
\tag{4.5.2}
\end{equation}
其中 $A_1$ 为 $r$ 阶方阵，$A_2$ 为 $n-r$ 阶方阵.

\noindent\textbf{证明}\quad 类似定理 4.5.1 即可证明. $\square$

显然推论 4.5.1 还可进一步推广. 设 $V=V_1\oplus V_2\oplus\cdots\oplus V_m$，其中每个 $V_i$ 都是线性变换 $\varphi$ 的不变子空间，那么在 $V$ 中存在一组基 (这组基可由 $V_i$ 的基合并而成)，使 $\varphi$ 在这组基下的表示矩阵为分块对角阵：
\[
\left(\begin{array}{cccc}
A_1 & & &\\
& A_2 & &\\
& & \ddots &\\
& & & A_m
\end{array}\right),
\]
其中 $A_i$ 是 $\varphi|_{V_i}$ 的表示矩阵，它是 $r_i$ 阶方阵，$r_i=\dim V_i$.

如果 $n$ 维线性空间的线性变换 $\varphi$ 有足够小的不变子空间，比如有 $n$ 个一维不变子空间，其直和正好组成全空间，那么上式中的表示矩阵就是一个对角阵.

\noindent\textbf{例 4.5.4}\quad 设 $V$ 是数域 $\mathbb{K}$ 上的三维空间，$\{e_1,e_2,e_3\}$ 是 $V$ 的基，$\varphi$ 是 $V$ 上的线性变换，它在这组基下的表示矩阵为
\[
\left(\begin{array}{rrr}
3 & 1 & -1\\
2 & 2 & -1\\
2 & 2 & 0
\end{array}\right).
\]
求证：由向量 $\{e_3,e_1+e_2+2e_3\}$ 生成的子空间 $U$ 是 $\varphi$ 的不变子空间.

\noindent\textbf{证明}\quad $\varphi(e_3)$ 的坐标向量为
\[
\left(\begin{array}{rrr}
3 & 1 & -1\\
2 & 2 & -1\\
2 & 2 & 0
\end{array}\right)
\left(\begin{array}{c}
0\\
0\\
1
\end{array}\right)
=
\left(\begin{array}{r}
-1\\
-1\\
0
\end{array}\right).
\]
容易求出向量组 $\{(-1,-1,0)',(0,0,1)',(1,1,2)'\}$ 的秩为 2，而 $U$ 显然是二维子空间，因此 $\varphi(e_3)\in U$. 同理可证 $\varphi(e_1+e_2+2e_3)\in U$，故 $U$ 是 $\varphi$-不变子空间. $\square$

\subsection*{习题 4.5}

1. 设线性空间 $V$ 上的线性变换 $\varphi$ 在一组基 $\{e_1,e_2,e_3,e_4\}$ 下的表示矩阵为
\[
A=
\left(\begin{array}{rrrr}
1 & 0 & 2 & -1\\
0 & 1 & 4 & -2\\
2 & -1 & 0 & 1\\
2 & -1 & -1 & 2
\end{array}\right),
\]
求证：$U=L(e_1+2e_2,e_3+e_4,e_1+e_2)$ 和 $W=L(e_2+e_3+2e_4)$ 都是 $\varphi$ 的不变子空间.

2. 设 $V_1,V_2$ 是 $V$ 上线性变换 $\varphi$ 的不变子空间，求证：$V_1+V_2,V_1\cap V_2$ 也是 $\varphi$ 的不变子空间.

3. 设 $\varphi,\psi$ 都是线性空间 $V$ 上的线性变换且 $\varphi\psi=\psi\varphi$，求证：$\operatorname{Im}\varphi$ 及 $\operatorname{Ker}\varphi$ 都是 $\psi$ 的不变子空间.

4. 设 $\varphi$ 是 $n$ 维线性空间 $V$ 上的自同构，若 $W$ 是 $\varphi$ 的不变子空间，求证：$W$ 也是 $\varphi^{-1}$ 的不变子空间.

5. 设 $V$ 是次数小于 $n$ 的实系数多项式组成的线性空间，$D$ 是 $V$ 上的求导变换，求证：

(1) $D$ 的任一 $r(1\leq r\leq n)$ 维不变子空间必由 $\{1,x,\cdots,x^{r-1}\}$ 生成的子空间；

(2) $\operatorname{Im}D\cap\operatorname{Ker}D\neq 0$ 且 $V\neq\operatorname{Im}D+\operatorname{Ker}D$.

6. 设 $\varphi$ 是 $n$ 维线性空间 $V$ 上的线性变换，$\varphi$ 在 $V$ 的一组基下的表示矩阵为对角阵且主对角线上的元素互不相同，求 $\varphi$ 的所有不变子空间.

\section*{历史与展望}

我们知道在线性同构的意义下 (参考定理 4.3.1)，抽象的线性空间之间的线性映射可以等同于具体的行 (列) 向量空间之间的线性映射，而后者其实就是坐标之间的线性变换.

Gauss 在著作《算术研究》中研究了二元二次型的算术理论. 若两个整系数二元二次型 $f(x,y)=ax^2+bxy+cy^2$ 与 $F(X,Y)=AX^2+BXY+CY^2$ 可诱导出相同的整数集，则称它们为等价的二次型. Gauss 证明了两个整系数二元二次型等价当且仅当存在坐标 $(x,y)$ 与 $(X,Y)$ 之间的线性变换 $T$ 将 $f(x,y)$ 变成 $F(X,Y)$，其中 $\det T=1$. 虽然 Gauss 并没有采用矩阵的术语，但他将线性变换表示为数的矩形阵列，并用坐标之间的线性变换的复合含蓄地定义了矩阵的乘积 (只对 $2\times 2$ 和 $3\times 3$ 的情形).

坐标之间的线性变换 $y_i=\sum_{j=1}^n a_{ij}x_j(1\leq i\leq m)$ 频繁出现在 17 世纪和 18 世纪解析几何的研究中 (主要是 $m=n\leq 3$ 的情形)，由此自然地诱导出矩阵的各种运算. 创立于 17 世纪的射影几何在 19 世纪早期用解析的语言进行描述，其中坐标之间的线性变换也开始出现.

Descartes 在 1637 年发表的著作《几何学》中用两个坐标来描述平面，这一工作不仅标志着解析几何的创立，开启了利用代数研究几何的先河，而且使“数形结合”的思想深入人心，深刻影响着数学的发展.

作为现代数学基石之一的高等代数，其实也是一门代数与几何交叉融合的学科. 行列式与矩阵是代数概念，线性空间与线性映射是几何概念. 在第四章中，我们建立了代数与几何之间的一座桥梁，实现了代数语言与几何语言之间的自由转换，从而可以把线性映射问题转化为矩阵问题，然后用代数方法来解决；也可以把矩阵问题转换为线性映射问题，然后用几何方法来解决.

实践证明：深入理解几何意义，熟练掌握代数方法，注重代数与几何之间的相互转换与有机统一，这就是学好高等代数的有效方法.

\section*{复习题四}

1. 设 $\varphi$ 是数域 $\mathbb{F}$ 上有限维线性空间 $V$ 到 $U$ 的线性映射，求证：必存在 $U$ 到 $V$ 的线性映射 $\psi$，使 $\varphi\psi\varphi=\varphi$.

2. 设有数域 $\mathbb{F}$ 上有限维线性空间 $V,V'$，又 $U$ 是 $V$ 的子空间，$\varphi$ 是 $U$ 到 $V'$ 的线性映射. 求证：必存在 $V$ 到 $V'$ 的线性映射 $\psi$，它在 $U$ 上的限制就是 $\varphi$ (称这样的 $\psi$ 是 $\varphi$ 的扩张).

3. 设 $V,U$ 是 $\mathbb{F}$ 上的有限维线性空间，$\varphi$ 是 $V$ 到 $U$ 的线性映射，求证：

(1) $\varphi$ 是单映射的充分必要条件是存在 $U$ 到 $V$ 的线性映射 $\psi$，使 $\psi\varphi=I_V$，这里 $I_V$ 表示 $V$ 上的恒等映射；

(2) $\varphi$ 是满映射的充分必要条件是存在 $U$ 到 $V$ 的线性映射 $\eta$，使 $\varphi\eta=I_U$，这里 $I_U$ 表示 $U$ 上的恒等映射.

4. 设 $V,U$ 是数域 $\mathbb{K}$ 上的有限维线性空间，$\varphi,\psi:V\to U$ 是两个线性映射，证明：存在 $U$ 上的线性变换 $\xi$，使得 $\psi=\xi\varphi$ 成立的充分必要条件是 $\operatorname{Ker}\varphi\subseteq\operatorname{Ker}\psi$.

5. 设 $V,U$ 是数域 $\mathbb{K}$ 上的有限维线性空间，$\varphi,\psi:V\to U$ 是两个线性映射，证明：存在 $V$ 上的线性变换 $\xi$，使得 $\psi=\varphi\xi$ 成立的充分必要条件是 $\operatorname{Im}\psi\subseteq\operatorname{Im}\varphi$.

6. 设 $V$ 是数域 $\mathbb{K}$ 上的 $n$ 维线性空间，$\varphi$ 是 $V$ 上的幂零线性变换，满足 $\operatorname{r}(\varphi)=n-1$. 求证：存在 $V$ 的一组基，使得 $\varphi$ 在这组基下的表示矩阵为
\[
\left(\begin{array}{ccccc}
0&0&\cdots&0&0\\
1&0&\cdots&0&0\\
0&1&\cdots&0&0\\
\vdots&\vdots&&\vdots&\vdots\\
0&0&\cdots&1&0
\end{array}\right).
\]

7. 设 $a_0,a_1,\cdots,a_n$ 是数域 $\mathbb{F}$ 上 $n+1$ 个不同的数，$V$ 是 $\mathbb{F}$ 上次数不超过 $n$ 的多项式全体组成的线性空间. 设 $\varphi$ 是 $V$ 到 $n+1$ 维行向量空间 $\mathbb{F}^{n+1}$ 的映射：
\[
\varphi(f)=(f(a_0),f(a_1),\cdots,f(a_n)),
\]
求证：$\varphi$ 是线性同构.

8. 设 $U_1,U_2$ 是 $n$ 维线性空间 $V$ 的子空间，假设它们维数相同，求证：存在 $V$ 上的可逆线性变换 $\varphi$，使 $U_2=\varphi(U_1)$.

9. 设 $\varphi$ 是 $n$ 维线性空间 $V$ 上的线性变换，若对 $V$ 中任一向量 $\alpha$，总存在正整数 $m$ ($m$ 可能和 $\alpha$ 有关)，使 $\varphi^m(\alpha)=0$. 求证：$I_V-\varphi$ 是自同构.

10. 设 $V=M_n(\mathbb{K})$，$A,B$ 是两个 $n$ 阶矩阵，定义 $V$ 上的变换：$\varphi(X)=AXB$. 求证：$\varphi$ 是 $V$ 上的线性变换，$\varphi$ 是可逆变换的充分必要条件是 $A,B$ 都是可逆阵.

11. 设 $\varphi$ 是 $n$ 维线性空间 $V$ 上的线性变换，若 $\varphi$ 是满映射，$U$ 是 $V$ 的 $r$ 维子空间，求证：$\varphi(U)$ 也是 $V$ 的 $r$ 维子空间. 举例说明：若 $V$ 是无限维线性空间，上述结论不成立.

12. 设 $V$ 是数域 $\mathbb{K}$ 上的线性空间 (不要求是有限维的)，$\varphi,\psi$ 是 $V$ 上的线性变换.

(1) 证明：$\varphi$ 和 $\psi$ 都是可逆变换的充分必要条件是 $\varphi\psi$ 和 $\psi\varphi$ 都是可逆变换.

(2) 若 $\varphi\psi=I_V$，则称 $\varphi$ 是 $\psi$ 的左逆变换，$\psi$ 是 $\varphi$ 的右逆变换. 证明：$\varphi$ 是可逆变换的充分必要条件是 $\varphi$ 有且仅有一个左逆变换 (右逆变换).

13. 试构造无限维线性空间 $V$ 以及 $V$ 上的线性变换 $\varphi,\psi$，使 $\varphi\psi-\psi\varphi=I_V$.

14. 设 $\varphi$ 是有限维线性空间 $V$ 到 $U$ 的线性映射，求证：必存在 $V$ 和 $U$ 的两组基，使线性映射 $\varphi$ 在这两组基下的表示矩阵为
\[
\left(\begin{array}{cc}
I_r & O\\
O & O
\end{array}\right).
\]

15. 设 $\varphi:V\to U$ 为线性映射且 $\varphi$ 的秩为 $r$，证明：存在 $r$ 个秩为 1 的线性映射 $\varphi_i:V\to U(i=1,\cdots,r)$，使 $\varphi=\varphi_1+\cdots+\varphi_r$.

16. 设 $\varphi$ 是线性空间 $V$ 上的线性变换，若它在 $V$ 的任一组基下的表示矩阵都相同，求证：$\varphi$ 是纯量变换，即存在常数 $k$，使 $\varphi(\alpha)=k\alpha$ 对一切 $\alpha\in V$ 都成立.

17. 设 $A,B$ 都是数域 $\mathbb{F}$ 上的 $m\times n$ 矩阵，求证：方程组 $Ax=0,Bx=0$ 同解的充分必要条件是存在 $m$ 阶非异阵 $P$，使 $B=PA$.

18. 设 $V$ 是数域 $\mathbb{F}$ 上 $n$ 阶矩阵全体构成的线性空间，$\varphi$ 是 $V$ 上的线性变换：$\varphi(A)=A'$. 证明：存在 $V$ 的一组基，使 $\varphi$ 在这组基下的表示矩阵是一个对角阵且主对角元素全是 1 或 $-1$，并求出 1 和 $-1$ 的个数.

19. 设 $V$ 是数域 $\mathbb{K}$ 上的 $n$ 维线性空间，$\varphi,\psi$ 是 $V$ 上的线性变换且 $\varphi^2=0,\psi^2=0,\varphi\psi+\psi\varphi=I$，$I$ 是 $V$ 上的恒等变换. 求证：

(1) $V=\operatorname{Ker}\varphi\oplus\operatorname{Ker}\psi$;

(2) 若 $V$ 是二维空间，则存在 $V$ 的基 $e_1,e_2$，使 $\varphi,\psi$ 在这组基下的表示矩阵分别为
\[
A=\left(\begin{array}{cc}
0&0\\
1&0
\end{array}\right),\quad
B=\left(\begin{array}{cc}
0&1\\
0&0
\end{array}\right);
\]

(3) $V$ 必是偶数维空间且若 $V$ 是 $2k$ 维空间，则存在 $V$ 的一组基，使 $\varphi,\psi$ 在这组基下的表示矩阵分别为下列分块对角阵：
\[
\left(\begin{array}{cccc}
A&O&\cdots&O\\
O&A&\cdots&O\\
\vdots&\vdots&&\vdots\\
O&O&\cdots&A
\end{array}\right),\quad
\left(\begin{array}{cccc}
B&O&\cdots&O\\
O&B&\cdots&O\\
\vdots&\vdots&&\vdots\\
O&O&\cdots&B
\end{array}\right),
\]
其中主对角线上分别有 $k$ 个 $A$ 和 $k$ 个 $B$.

20. 设 $V$ 是数域 $\mathbb{F}$ 上的线性空间，$\varphi_1,\varphi_2,\cdots,\varphi_k$ 是 $V$ 上的非零线性变换. 求证：存在 $\alpha\in V$，使 $\varphi_i(\alpha)\neq 0(i=1,2,\cdots,k)$.

21. 设 $V$ 是数域 $\mathbb{F}$ 上的线性空间，$\varphi_1,\varphi_2,\cdots,\varphi_k$ 是 $V$ 上互不相同的线性变换. 求证：存在 $\alpha\in V$，使 $\varphi_1(\alpha),\varphi_2(\alpha),\cdots,\varphi_k(\alpha)$ 互不相同.

22. 设 $A$ 是 $n$ 阶方阵，求证：$\operatorname{r}(A^n)=\operatorname{r}(A^{n+1})=\operatorname{r}(A^{n+2})=\cdots$.

23. 设 $\varphi$ 是 $n$ 维线性空间 $V$ 上的线性变换，求证：必存在整数 $m\in[0,n]$，使
\[
\operatorname{Im}\varphi^m=\operatorname{Im}\varphi^{m+1},\quad \operatorname{Ker}\varphi^m=\operatorname{Ker}\varphi^{m+1},\quad V=\operatorname{Im}\varphi^m\oplus\operatorname{Ker}\varphi^m.
\]

24. 设 $V$ 是数域 $\mathbb{K}$ 上的 $n$ 维线性空间，$\varphi$ 是 $V$ 上的线性变换，证明以下 9 个结论等价：

(1) $V=\operatorname{Ker}\varphi\oplus\operatorname{Im}\varphi$;

(2) $V=\operatorname{Ker}\varphi+\operatorname{Im}\varphi$;

(3) $\operatorname{Ker}\varphi\cap\operatorname{Im}\varphi=0$;

(4) $\operatorname{Ker}\varphi=\operatorname{Ker}\varphi^2$，或等价地，$\dim\operatorname{Ker}\varphi=\dim\operatorname{Ker}\varphi^2$;

(5) $\operatorname{Ker}\varphi=\operatorname{Ker}\varphi^2=\operatorname{Ker}\varphi^3=\cdots$，或等价地，$\dim\operatorname{Ker}\varphi=\dim\operatorname{Ker}\varphi^2=\dim\operatorname{Ker}\varphi^3=\cdots$;

(6) $\operatorname{Im}\varphi=\operatorname{Im}\varphi^2$，或等价地，$\operatorname{r}(\varphi)=\operatorname{r}(\varphi^2)$;

(7) $\operatorname{Im}\varphi=\operatorname{Im}\varphi^2=\operatorname{Im}\varphi^3=\cdots$，或等价地，$\operatorname{r}(\varphi)=\operatorname{r}(\varphi^2)=\operatorname{r}(\varphi^3)=\cdots$;

(8) $\operatorname{Ker}\varphi$ 存在 $\varphi$-不变补空间，即存在 $\varphi$-不变子空间 $U$，使得 $V=\operatorname{Ker}\varphi\oplus U$;

(9) $\operatorname{Im}\varphi$ 存在 $\varphi$-不变补空间，即存在 $\varphi$-不变子空间 $W$，使得 $V=\operatorname{Im}\varphi\oplus W$.

25. 设 $V,U$ 是数域 $\mathbb{K}$ 上的有限维线性空间，$\varphi:V\to U$ 是线性映射，证明：由 $\varphi$ 诱导的线性映射 $\overline{\varphi}:V/\operatorname{Ker}\varphi\to\operatorname{Im}\varphi$，$\overline{\varphi}(v+\operatorname{Ker}\varphi)=\varphi(v)$ 是线性同构. 特别地，$\dim V=\dim\operatorname{Ker}\varphi+\dim\operatorname{Im}\varphi$.

26. 设 $U,W$ 是 $n$ 维线性空间 $V$ 的子空间且 $\dim U+\dim W=\dim V$. 求证：存在 $V$ 上的线性变换 $\varphi$，使 $\operatorname{Ker}\varphi=U,\operatorname{Im}\varphi=W$.

27. 设 $\varphi$ 是有限维线性空间 $V$ 到 $U$ 的满线性映射，求证：必存在 $V$ 的子空间 $W$，使 $V=W\oplus\operatorname{Ker}\varphi$，且 $\varphi$ 在 $W$ 上的限制是 $W$ 到 $U$ 上的线性同构.

28. 设 $\varphi$ 是 $n(n\geq 2)$ 维线性空间 $V$ 上的线性变换，证明以下 $n$ 个结论等价：

(1) $V$ 的任一 1 维子空间都是 $\varphi$-不变子空间；

$\cdots\cdots\cdots$

(r) $V$ 的任一 $r$ 维子空间都是 $\varphi$-不变子空间；

$\cdots\cdots\cdots$

(n-1) $V$ 的任一 $n-1$ 维子空间都是 $\varphi$-不变子空间；

(n) $\varphi$ 是纯量变换.

29. 设 $A$ 为数域 $\mathbb{K}$ 上的 $n$ 阶幂零阵，$B$ 为 $n$ 阶方阵，满足 $AB=BA$ 且 $\operatorname{r}(AB)=\operatorname{r}(B)$. 求证：$B=O$.

30. 设 $\varphi$ 是 $n$ 维线性空间 $V$ 上的线性变换，$U$ 是 $r$ 维 $\varphi$-不变子空间. 取 $U$ 的一组基 $\{e_1,\cdots,e_r\}$，并扩张为 $V$ 的一组基 $\{e_1,\cdots,e_r,e_{r+1},\cdots,e_n\}$. 设 $\varphi$ 在这组基下的表示矩阵
\[
A=(a_{ij})=
\left(\begin{array}{cc}
A_{11} & A_{12}\\
O & A_{22}
\end{array}\right)
\]
为分块上三角阵，其中 $A_{11}$ 是 $\varphi$ 在不变子空间 $U$ 上的限制 $\varphi|_U$ 在基 $\{e_1,\cdots,e_r\}$ 下的表示矩阵. 证明：$\varphi$ 诱导的变换 $\overline{\varphi}(v+U)=\varphi(v)+U$ 是商空间 $V/U$ 上的线性变换，并且在 $V/U$ 的一组基 $\{e_{r+1}+U,\cdots,e_n+U\}$ 下的表示矩阵为 $A_{22}$.

31. 设 $\varphi$ 是 $n$ 维线性空间 $V$ 上的幂等线性变换，即满足 $\varphi^2=\varphi$. 证明：$V=U\oplus W$，其中 $U=\operatorname{Im}\varphi=\operatorname{Ker}(I_V-\varphi)$，$W=\operatorname{Im}(I_V-\varphi)=\operatorname{Ker}\varphi$.

32. 设 $A$ 是数域 $\mathbb{F}$ 上的 $n$ 阶幂等阵，求证：

(1) 存在 $n$ 阶非异阵 $P$，使
\[
P^{-1}AP=
\left(\begin{array}{cc}
I_r & O\\
O & O
\end{array}\right),
\]
其中 $r=\operatorname{r}(A)$;

(2) $\operatorname{r}(A)=\operatorname{tr}(A)$.

33. 设 $A,B$ 是数域 $\mathbb{F}$ 上的 $n$ 阶幂等阵，且 $A$ 和 $B$ 的秩相同，求证：必存在 $\mathbb{F}$ 上的 $n$ 阶可逆阵 $C$，使 $CB=AC$.

34. 设 $\varphi,\psi$ 是 $n$ 维线性空间 $V$ 上的幂等线性变换，求证：

(1) $\operatorname{Im}\varphi=\operatorname{Im}\psi$ 的充分必要条件是 $\varphi\psi=\psi,\ \psi\varphi=\varphi$;

(2) $\operatorname{Ker}\varphi=\operatorname{Ker}\psi$ 的充分必要条件是 $\varphi\psi=\varphi,\ \psi\varphi=\psi$;

(3) $\varphi+\psi$ 是幂等变换的充分必要条件是 $\varphi\psi=\psi\varphi=0$;

(4) $\varphi-\psi$ 是幂等变换的充分必要条件是 $\varphi\psi=\psi\varphi=\psi$.

35. 设 $\varphi_1,\cdots,\varphi_m$ 是 $n$ 维线性空间 $V$ 上的线性变换，且适合条件：
\[
\varphi_i^2=\varphi_i,\quad \varphi_i\varphi_j=0(i\neq j),\quad \operatorname{Ker}\varphi_1\cap\cdots\cap\operatorname{Ker}\varphi_m=0.
\]
求证：$V$ 是 $\operatorname{Im}\varphi_1,\cdots,\operatorname{Im}\varphi_m$ 的直和.

36. 设 $\varphi,\varphi_1,\cdots,\varphi_m$ 是 $n$ 维线性空间 $V$ 上的线性变换，满足：$\varphi^2=\varphi$ 且 $\varphi=\varphi_1+\varphi_2+\cdots+\varphi_m$. 求证：$\operatorname{r}(\varphi)=\operatorname{r}(\varphi_1)+\operatorname{r}(\varphi_2)+\cdots+\operatorname{r}(\varphi_m)$ 成立的充分必要条件是 $\varphi_i^2=\varphi_i,\ \varphi_i\varphi_j=0(i\neq j)$.

\end{document}
